% Generated by GrindEQ Word-to-LaTeX 
\documentclass{article} % use \documentstyle for old LaTeX compilers



\usepackage[utf8]{inputenc} % 'cp1252'-Western, 'cp1251'-Cyrillic, etc.
\usepackage[english]{babel} % 'french', 'german', 'spanish', 'danish', etc.
\usepackage{amsmath}
\usepackage{amssymb}
\usepackage{txfonts}
\usepackage{mathdots}
\usepackage[classicReIm]{kpfonts}
\usepackage{graphicx}



% You can include more LaTeX packages here 





\begin{document}



%\selectlanguage{english} % remove comment delimiter ('%') and select language if required





\noindent 



\noindent 



\noindent 



\noindent 



\noindent \textbf{Grey Relational Analysis: A Statistical Approach to Measuring Dependence in Data-Poor Conditions}



\noindent Mohamad Ahmadian ${}^{a}$ \footnote{ $\ Corresponding\ author.$  $E-mail\ addresses:\ ahmadian333@ut.ac.ir$ }${}^{\ }$



\noindent ${}^{a\ }$Faculty of Management and Accounting, College of Farabi, University of Tehran, Tehran, Iran (ahmadian333@ut.ac.ir)



\noindent \textbf{Abstract}



\noindent Grey relational analysis is one of the most important tools in grey systems theory that is used to measure the intensity of the relationship between variables in situations where data are limited, incomplete, or uncertain. In terms of functionality, this method is similar to some statistical methods such as correlation and dependence analysis, but unlike classical statistics, it does not depend on strict assumptions such as normality of distribution, large sample size, or linearity of the relationship. Grey relational analysis, by focusing on the similarity of the pattern of changes between data series, allows ranking of effective factors and identifying important variables. In this article, the theoretical foundations, formulation, implementation steps, its connection with statistical concepts, and its applications in data analysis are reviewed.



\noindent \textbf{Keywords}: Grey relational analysis, grey statistics, statistical dependence, correlation, incomplete data, factor ranking



\noindent 



\begin{enumerate}
\item  \textbf{Introduction}
\end{enumerate}



\noindent In many statistical studies, the main goal is to examine the relationship between variables. Researchers usually want to know which factor has the greatest impact on an indicator, which variables are dependent on each other, and to what extent this dependence is. In classical statistics, tools such as Pearson's correlation coefficient, Spearman's correlation, regression, and analysis of variance are used to answer these questions. However, these methods often depend on conditions such as a sufficient number of data, appropriate quality of observations, and the establishment of some statistical assumptions (such as normality of distribution).



\noindent In real problems, especially in the humanities, management, economics, environmental science, and some engineering fields, data are often limited and incomplete. In such situations, Grey Relational Analysis (GRA) is proposed as a complementary and, in some cases, alternative method for measuring the dependence between variables. This method, which belongs to the field of ``Grey Systems Theory'', is designed for systems with incomplete information or high uncertainty (Ahmadian, 2026).



\noindent Unlike classical methods based on probability distributions, GRA is based on the ``geometric similarity'' between data sequences. Simply put, this method examines whether the trend of changes in one variable (reference sequence) is consistent with the trend of changes in other variables (comparison sequences). The closer the shape of the data curves, the stronger the grey relationship and the higher the grey correlation coefficient (Shi et al., 2026).



\noindent The main advantage of this method is the lack of strict statistical assumptions and high efficiency with small data volumes. This feature has made GRA an ideal tool for decision-making in complex and uncertain situations. In the following, new applications of this method in scientific articles are reviewed\textbf{.}



\begin{enumerate}
\item \textbf{ Basics of Grey Relational Analysis and its Relationship to Statistics}
\end{enumerate}



\noindent This section provides a precise definition of gray relational analysis and its place in statistics. It then examines how this method fits into the framework of statistical analysis and what role it plays in complementing classical methods.



\noindent \textit{2.1 Definition and Concept of Grey Relational Analysis}



\noindent Grey relational analysis is a method for measuring the degree of closeness or similarity between a reference sequence and several comparison sequences (Liu et al., 2022). If the changes of a variable during observations have a pattern similar to the changes of the target variable, it is said that the degree of Grey relational of that variable with the target variable is high. As a result, that variable is considered a more important factor in explaining the behavior of the dependent variable from a statistical and analytical point of view.



\noindent From a statistical perspective, Grey relational analysis is a type of non-parametric dependence measurement. In this method, instead of estimating distribution parameters or calculating covariance, the relative distance between normalized data is used. Therefore, this method is very suitable for data that: have a small sample size, do not follow a specific distribution, have different scales, have irregularity and incomplete information.



\noindent \textit{2.2 The Place of Grey Relational Analysis in Statistics}



\noindent Grey relational analysis is similar in purpose to statistical methods of measuring relationships, but it differs from them in basis. In statistics, correlation indicates the degree of covariation between two variables. For example, Pearson's correlation coefficient focuses on the linear relationship and numerical structure of the data. In contrast, Grey relational analysis focuses on the geometric similarity and pattern of changes in sequences.



\noindent In statistical terms: If correlation is concerned with ``quantitative covariance'', Grey relational analysis is concerned with ``co-direction and co-formity of trends''. Therefore, two variables may not have a high correlation coefficient, but they are close to each other in terms of Grey relational analysis, because their pattern of increase and decrease is similar. This feature has made Grey relational analysis a useful tool for exploratory data analysis, variable selection, and factor ranking.



\noindent Grey relational analysis is often compared to the correlation coefficient, but there are important differences between them.



\noindent First, Pearson correlation measures the strength of a linear relationship between two variables, while the Grey relational analysis does not depend on linearity. Second, classical correlation usually requires a larger sample size for more robust results, while the Grey relational analysis can also be applied to small samples. Third, in correlation, the presence of outliers can have a large effect on the result, but in Grey relational analysis, due to the emphasis on the pattern of change, this effect is in some cases less. However, it should be noted that the Grey relational analysis is not a complete substitute for correlation. Correlation is an accurate statistical indicator for measuring dependence in a probability framework, but the Grey relational analysis is more of a tool for analyzing poorly informed systems and ranking factors in conditions of insufficient data.



\begin{enumerate}
\item  \textbf{Grey Relational Analysis Method}
\end{enumerate}



\noindent The calculation method for GRA is as follows (Tharisanan et al., 2026):



\noindent Suppose a reference sequence is available as Eq.1.



\begin{tabular}{|p{0.9in}|p{2.6in}|p{0.9in}|} \hline 
 & $X_{0} =\{ x_{0} \eqref{GrindEQ__1_},x_{0} \eqref{GrindEQ__2_},...,x_{0} (n)\} $ &   \\ \hline 
\end{tabular}



\textit{m} comparison sequences are also defined as Eq.2.



\begin{tabular}{|p{0.9in}|p{2.6in}|p{0.9in}|} \hline 
 & $X_{i} =\{ x_{i} \eqref{GrindEQ__1_},x_{i} \eqref{GrindEQ__2_},\mathrm{...},x_{i} (n)\} ,\quad i=1,2,\mathrm{...},m$ & \eqref{GrindEQ__2_} \\ \hline 
\end{tabular}



The goal is to calculate the degree of relational of each comparison sequence to the reference sequence.



\noindent Since the variables may be measured on different scales, the data must first be normalized. This step is very important in statistics, because direct comparison of dissimilar data can be misleading. Depending on the nature of the indicator, different methods of scale-independence are used; such as: initial normalization, conversion to mean ratio, and scaling between zero and one. Normalization allows the analysis to focus on the pattern of changes, rather than the absolute magnitude of the numbers.



\noindent After normalization, the absolute difference between the reference sequence and each comparison sequence is calculated for each observation is available as Eq.3.



\begin{tabular}{|p{0.9in}|p{2.6in}|p{0.9in}|} \hline 
 & $\mathrm{\Delta }_{i} (k)=\mathrm{}x_{0} (k)-x_{i} (k)\mathrm{}$ & \eqref{GrindEQ__3_} \\ \hline 
\end{tabular}



This difference is statistically a measure of the distance between corresponding observations.



\noindent Then, in the entire data, the smallest and largest difference values are determined as Eq.4 and Eq.5.



\begin{tabular}{|p{0.9in}|p{2.6in}|p{0.9in}|} \hline 
 & $\mathrm{\Delta }_{\min } ={\mathop{\min }\limits_{i}} {\mathop{\min }\limits_{k}} \mathrm{\Delta }_{i} (k)$ & \eqref{GrindEQ__4_} \\ \hline 
 & $\mathrm{\Delta }_{\max } ={\mathop{\max }\limits_{i}} {\mathop{\max }\limits_{k}} \mathrm{\Delta }_{i} (k)$ & \eqref{GrindEQ__5_} \\ \hline 
\end{tabular}



These two values are used to scale the distances and calculate the relational coefficient.



\noindent The Grey relational coefficient for sequence i at point k is obtained from Eq.6.



\begin{tabular}{|p{0.9in}|p{2.6in}|p{0.9in}|} \hline 
 & $\xi _{i} (k)=\frac{\mathrm{\Delta }_{\min } +\rho \mathrm{\Delta }_{\max } }{\mathrm{\Delta }_{i} (k)+\rho \mathrm{\Delta }_{\max } } $ & \eqref{GrindEQ__6_} \\ \hline 
\end{tabular}



In this relationship, $\rhoup$ is the discrimination coefficient and its value is usually taken to be 0.5. This parameter plays the role of regulating the sensitivity of the model from a statistical point of view. The smaller $\rhoup$ is, the greater the discrimination between the sequences.



\noindent The value of $\xiup$i(k) is between zero and one. The closer this value is to one, the greater the similarity of that observation from the comparison sequence to the reference sequence.



\noindent To obtain an overall index of the degree of dependence of each variable with the reference variable, the average of the relational coefficients is calculated in Eq.7.



\begin{tabular}{|p{0.9in}|p{2.6in}|p{0.9in}|} \hline 
 & $\gamma _{i} =\frac{1}{n} \mathop{\sum }\limits_{k=1}^{n} \xi _{i} (k)$ & \eqref{GrindEQ__7_} \\ \hline 
\end{tabular}



The Grey relational degree $\gammaup$i is actually a summary statistic that shows the intensity of the overall relationship between the two sequences. The larger this value, the stronger the statistical relationship between the variable of interest and the reference variable is considered.



\noindent 



\begin{enumerate}
\item  \textbf{Example of Grey Relational Analysis}
\end{enumerate}



\noindent To illustrate the application of grey relational analysis, consider a bank that seeks to identify which operational indicator is most closely associated with the monthly profit of a branch. Three explanatory variables are considered: the number of active customers (X${}_{1}$), the number of digital transactions (X${}_{2}$), and the number of customer complaints (X${}_{3}$). The response variable (X${}_{0}$) is the monthly profit (thousand USD). Table 1 presents a simple illustrative dataset.



\noindent \textit{Table 1 Sample data}



\begin{tabular}{|p{0.7in}|p{1.0in}|p{0.9in}|p{0.9in}|p{0.9in}|} \hline 
Month & X${}_{0}$ & X${}_{1}$ & X${}_{2}$ & X${}_{3}$ \\ \hline 
1 & 100 & 96 & 99 & 80 \\ \hline 
2 & 120 & 115 & 120 & 95 \\ \hline 
3 & 140 & 133 & 141 & 120 \\ \hline 
\end{tabular}



The absolute differences are shown in Table 2.



\noindent \textit{Table 2 Absolute difference sequences}



\begin{tabular}{|p{0.7in}|p{1.3in}|p{1.2in}|p{1.2in}|} \hline 
Month & $\Delta$1 & $\Delta$2 & $\Delta$3 \\ \hline 
1 & 4 & 1 & 20 \\ \hline 
2 & 5 & 0 & 25 \\ \hline 
3 & 7 & 1 & 20 \\ \hline 
\end{tabular}



Thus, $\mathrm{\Delta }_{\min } =0$, $\mathrm{\Delta }_{\max } =25$



\noindent The grey relational coefficients, calculated using the commonly adopted distinguishing coefficient $\rho $=0.5, are presented in Table 3.



\noindent \textit{Table 3 Grey relational coefficients}



\begin{tabular}{|p{0.7in}|p{1.3in}|p{1.2in}|p{1.2in}|} \hline 
\textbf{Month} & \textbf{$\boldsymbol{\xiup}$1} & \textbf{$\boldsymbol{\xiup}$2{}} & \textbf{$\boldsymbol{\xiup}$3{}} \\ \hline 
\textbf{1} & \textbf{0.758} & \textbf{0.926} & \textbf{0.385} \\ \hline 
\textbf{2} & \textbf{0.714} & \textbf{1.000} & \textbf{0.333} \\ \hline 
\textbf{3} & \textbf{0.641} & \textbf{0.926} & \textbf{0.385} \\ \hline 
\end{tabular}



The grey relational grades are obtained by averaging the coefficients. Accordingly, $\gammaup$1{}=0.704, $\gammaup$2=0.951, $\gammaup$3=0.368, hence, $\gammaup$2{}$\mathrm{>}$$\gammaup$1{}$\mathrm{>}$$\gammaup$3.



\noindent The results indicate that X${}_{2\ }$(digital transactions) exhibit the strongest grey relational grade with monthly profit, followed by the number of active customers, whereas customer complaints have the weakest relationship. This outcome is consistent with managerial expectations, as higher digital transaction volumes are generally associated with increased banking activity and revenue generation, while customer complaints are less directly related to profitability. Consequently, grey relational analysis provides a practical tool for ranking explanatory variables according to their association with a response variable, thereby supporting feature selection, dimensionality reduction, and sensitivity analysis.



\begin{enumerate}
\item  \textbf{Applications, Advantages, and Limitations}
\end{enumerate}



\noindent In this section, the practical applications of gray relational analysis in various domains are reviewed. The advantages and limitations of this method are also analyzed in detail to provide a complete view of its capabilities and challenges.



\noindent \textit{5.1 Applications}



\noindent This method is used in various fields for data analysis. In economics, it can be used to determine the factors affecting economic growth or inflation. In management, it is used to rank organizational performance indicators. In social sciences, it is used to measure the proximity of variables such as satisfaction, development, or quality of life. In engineering and environmental science, it is also used to evaluate parameters affecting quality, sustainability, or pollution.



\noindent In all these cases, Grey relational analysis plays a role similar to a statistical tool to discover the structure of dependence between variables. For example, articles can be reviewed in the following areas: Economic studies that have used Grey relational analysis to identify factors affecting economic growth, inflation, productivity, or financial markets. Management research that has used this method to evaluate organizational performance, rank productivity indicators, select suppliers, or analyze customer satisfaction. Environmental articles that have analyzed the relationship between pollution indicators, energy consumption, water quality, or sustainable development using this method. Engineering and industrial research that has used Grey relational analysis to select optimal parameters, quality control, or assess reliability. Health and social science studies in which this method has been used to rank risk factors, analyze well-being indicators, or examine variables affecting quality of life.



\noindent \textit{5.2 Advantages of Grey relational analysis in statistical research}



\noindent Grey relational analysis has several important advantages that make it attractive for applied statistical studies: Requirement of small sample size, no strong dependence on distributional assumptions, ability to use incomplete or irregular data, possibility of comparing non-scale variables after normalization, suitable for ranking effective factors, relative simplicity in calculation and interpretation. For this reason, this method is widely used in research where data are limited or conditions for performing classical tests are not available.



\noindent \textit{5.3 Limitations of the Method}



\noindent Despite its many advantages, Grey relational analysis also has some limitations. Its results may be sensitive to the normalization method. Also, the choice of the discrimination coefficient can affect the final ranking. On the other hand, this method mostly shows the intensity of similarity and, like regression models, does not provide an explicit causal or functional relationship.



\noindent As a result, from a statistical perspective, Grey relational analysis is better used as a complementary method alongside other methods such as correlation, regression, factor analysis, or multi-criteria decision-making methods.



\begin{enumerate}
\item  \textbf{Review of recent applied studies}
\end{enumerate}



\noindent In this section, recent research that has used GRA in diverse areas such as managing large projects, selecting sustainable suppliers in the electric vehicle supply chain, and accurately predicting fuel properties is reviewed. These studies show how combining GRA with other methods (such as fuzzy evaluation or machine learning) can be a powerful tool for decision-making under conditions of uncertainty and high complexity.



\noindent Anjum et al. (2026) used Grey relational analysis to rank AI-based mental health interventions. By collecting opinions from experts and patients, it was shown that ``therapeutic effectiveness'' was the most important criterion for selecting these tools, and this priority did not differ between different demographic groups. It was also shown that the grey model can prioritize complex criteria well when compared to other methods. Xiao et al. (2026) used a dynamic lag Grey relational analysis model to show that ``new productivity'' (such as Internet development and technological innovation) reduces rural energy poverty with a time lag. It was proved that if time and lag effects are considered, the grey model can show the relationship between technology and energy poverty reduction much more accurately than conventional methods, and its effect is greater in less developed areas. Jin (2026) combined fuzzy evaluation and Grey relational analysis to provide a model to quantitatively measure the complexity of large construction projects. Using the Grey method, the weight and importance of each complexity factor (such as technical, organizational, and environmental) were calculated, and it was shown that this method can provide the most accurate picture of the complexity status of large projects without subjective bias. Singh \& Pandey (2026) used Grey relational analysis to select the best supplier in the electric vehicle supply chain with a focus on environmental sustainability. They examined 17 environmental criteria and showed that the method could rank suppliers well, and the fourth supplier was identified as the best option due to its lower resource consumption and use of new technologies.  Lawal et al. (2025) used Grey relational analysis to build a hybrid machine learning model to improve the prediction of the calorific value of coal. It was shown that the individual models performed poorly when exposed to new data, but by combining them using the grey method, the prediction accuracy was significantly improved.



\begin{enumerate}
\item  \textbf{ Conclusion}
\end{enumerate}



\noindent Grey relational analysis is an efficient method for measuring the dependence and ranking of variables in situations where data are limited, incomplete, or lack the assumptions necessary for classical statistical analyses. By focusing on the similarity of the pattern of changes between sequences, this method provides a non-parametric statistical measure of the relationship between variables. Although it cannot be a complete replacement for classical methods such as correlation and regression, it is very useful in exploratory analyses, variable selection, factor evaluation, and decision-making. Hence, Grey relational analysis can be considered a bridge between applied statistics and grey systems theory.



\noindent \textbf{References}



\noindent Ahmadian, M. (2026). Forecasting inbound tourist arrivals to Iran post-COVID-19 pandemic: a small data approach using grey models. Journal of Tourism Futures, 1--15. doi:https://doi.org/10.1108/JTF-10-2024-0219Anjum, R., Daud, S., Bhatti, G., \& Saddique, F. (2026). Evaluation of AI-based mental health interventions using the grey relational analysis. Grey Systems: Theory and Application, 16\eqref{GrindEQ__1_}, 75--90. doi:https://doi.org/10.1108/GS-03-2025-0031Jin, S. (2026). Measuring complexity in mega construction projects: fuzzy comprehensive evaluation and grey relational analysis. Engineering, Construction and Architectural Management, 33\eqref{GrindEQ__1_}, 284--317. doi:https://doi.org/10.1108/ECAM-07-2024-0951Lawal, A., Ajeboriogbon, A., Onifade, M., Oluwaseyi Bada, S., \& Mulenga, F. (2026). A novel grey relational analysis-based committee of machine learning methods for enhanced prediction of coal calorific value. Fuel, 406 Part C, 137070. doi:https://doi.org/10.1016/j.fuel.2025.137070Liu, S., Yang, Y., \& Forrest, J.-L. (2022). Grey Relational Analysis Models. In Grey Systems Analysis (pp. 77--124). Singapore: Springer. doi:https://doi.org/10.1007/978-981-19-6160-1\_5Shi, L., Luo, D., Zhang, M., \& Zhang, D. (2026). Research on the correlation degree between the manufacturing industry and the digital industry in the Yellow River Basin based on the grey matrix correlation model Available to Purchase. Grey Systems: Theory and Application, 1-23. doi:https://doi.org/10.1108/GS-12-2025-0194Singh, G., \& Pandey, A. (2026). Environmental sustainability integrated supplier selection in electric vehicle supply chains: a grey relational analysis approach. Environment, Development and Sustainability, 28, 7861--7889. doi:https://doi.org/10.1007/s10668-024-05294-xTharisanan, R., Kathiresan, M., Kathiravan, K., \& Priyavarshini, M. (2026). Multiobjective optimization of machining parameters on electrical discharge machine by Grey Relational Analysis coupled with Firefly Algorithm. Journal of the Chinese Institute of Engineers, 49\eqref{GrindEQ__1_}. doi:https://doi.org/10.1080/02533839.2025.2544774Xiao, Q., Gao, M., \& Wu, R. (2026). How can new quality productivity alleviate rural energy poverty? Evidence from dynamic lag grey relational analysis. Energy Sources, Part B: Economics, Planning, and Policy, 21\eqref{GrindEQ__1_}. doi:https://doi.org/10.1080/15567249.2025.2599195



\noindent 



\noindent 





\end{document}



